\subsection*{The Extended Successor}

\begin{tcolorbox}[colback=defbox, colframe=defborder, arc=2pt,
  left=6pt, right=6pt, top=4pt, bottom=4pt,
  title={\small\textbf{Definition (Successor on \texorpdfstring{$\mathbb{W}$}{W})}},
  fonttitle=\small\bfseries]
\begin{definition}[Successor on $\mathbb{W}$]
\label{def:successor-on-w}
Define $S_{\mathbb{W}}:\mathbb{W}\to\mathbb{W}$ by
\[
S_{\mathbb{W}}(0)=1
\qquad\text{and}\qquad
S_{\mathbb{W}}(n)=S_{\mathbb{N}}(n)\quad(n\in\mathbb{N}).
\]
\end{definition}
\end{tcolorbox}

\begin{remark*}[Interpretation]
The definition inserts $0$ immediately before the positive chain. After the
first step, successor is exactly the successor already defined on
$\mathbb{N}$.
\end{remark*}

\begin{dependencies}
\begin{itemize}
  \item \hyperref[def:whole-numbers]{Whole Numbers}
  \item \hyperref[def:peano-system]{Peano System}
\end{itemize}
\end{dependencies}

\begin{tcolorbox}[colback=thmbox, colframe=thmborder, arc=2pt,
  left=6pt, right=6pt, top=4pt, bottom=4pt,
  title={\small\textbf{Theorem (Successor on \texorpdfstring{$\mathbb{W}$}{W} Is Well-Defined)}},
  fonttitle=\small\bfseries]
\begin{theorem}[Successor on $\mathbb{W}$ Is Well-Defined]
\label{thm:successor-on-w-well-defined}
The rule defining $S_{\mathbb{W}}$ determines a function
$S_{\mathbb{W}}:\mathbb{W}\to\mathbb{W}$.
\hyperref[prf:successor-on-w-well-defined]{Go to proof.}
\end{theorem}
\end{tcolorbox}

\begin{remark*}[Interpretation]
The cases are exclusive because $0\notin\mathbb{N}$, and the positive case
lands in $\mathbb{N}\subseteq\mathbb{W}$.
\end{remark*}

\begin{dependencies}
\begin{itemize}
  \item \hyperref[def:whole-numbers]{Whole Numbers}
  \item \hyperref[def:successor-on-w]{Successor on $\mathbb{W}$}
\end{itemize}
\end{dependencies}

\begin{tcolorbox}[colback=thmbox, colframe=thmborder, arc=2pt,
  left=6pt, right=6pt, top=4pt, bottom=4pt,
  title={\small\textbf{Theorem (\texorpdfstring{$0$}{0} Is Not a Successor in \texorpdfstring{$\mathbb{W}$}{W})}},
  fonttitle=\small\bfseries]
\begin{theorem}[$0$ Is Not a Successor in $\mathbb{W}$]
\label{thm:zero-is-not-successor-in-w}
There is no $a\in\mathbb{W}$ such that $S_{\mathbb{W}}(a)=0$.
\hyperref[prf:zero-is-not-successor-in-w]{Go to proof.}
\end{theorem}
\end{tcolorbox}

\begin{remark*}[Interpretation]
The new element $0$ is the first element of the extended chain. It is placed
before $1$, not obtained by applying successor to some earlier element.
\end{remark*}

\begin{dependencies}
\begin{itemize}
  \item \hyperref[def:successor-on-w]{Successor on $\mathbb{W}$}
  \item \hyperref[ax:peano-one-not-successor]{$1$ Is Not a Successor}
\end{itemize}
\end{dependencies}

\begin{tcolorbox}[colback=thmbox, colframe=thmborder, arc=2pt,
  left=6pt, right=6pt, top=4pt, bottom=4pt,
  title={\small\textbf{Theorem (Every Nonzero Element of \texorpdfstring{$\mathbb{W}$}{W} Is a Successor)}},
  fonttitle=\small\bfseries]
\begin{theorem}[Every Nonzero Element of $\mathbb{W}$ Is a Successor]
\label{thm:every-nonzero-w-is-successor}
If $a\in\mathbb{W}$ and $a\neq 0$, then there exists
$b\in\mathbb{W}$ such that $a=S_{\mathbb{W}}(b)$.
\hyperref[prf:every-nonzero-w-is-successor]{Go to proof.}
\end{theorem}
\end{tcolorbox}

\begin{remark*}[Interpretation]
The element $1$ is the successor of $0$, and every positive natural number
beyond $1$ remains a successor in the original Peano chain.
\end{remark*}

\begin{dependencies}
\begin{itemize}
  \item \hyperref[def:whole-numbers]{Whole Numbers}
  \item \hyperref[def:successor-on-w]{Successor on $\mathbb{W}$}
  \item \hyperref[ax:peano-induction]{Induction Axiom}
\end{itemize}
\end{dependencies}

\begin{tcolorbox}[colback=thmbox, colframe=thmborder, arc=2pt,
  left=6pt, right=6pt, top=4pt, bottom=4pt,
  title={\small\textbf{Theorem (Successor on \texorpdfstring{$\mathbb{W}$}{W} Is Injective)}},
  fonttitle=\small\bfseries]
\begin{theorem}[Successor on $\mathbb{W}$ Is Injective]
\label{thm:successor-on-w-injective}
For all $a,b\in\mathbb{W}$,
\[
S_{\mathbb{W}}(a)=S_{\mathbb{W}}(b)\Longrightarrow a=b.
\]
\hyperref[prf:successor-on-w-injective]{Go to proof.}
\end{theorem}
\end{tcolorbox}

\begin{remark*}[Interpretation]
The extended successor still never merges two inputs. The only new case is the
comparison between $0$ and a positive natural number, and that case is ruled
out by the original non-return condition for the first element.
\end{remark*}

\begin{dependencies}
\begin{itemize}
  \item \hyperref[def:successor-on-w]{Successor on $\mathbb{W}$}
  \item \hyperref[ax:peano-successor-injective]{Injectivity of Successor}
  \item \hyperref[ax:peano-one-not-successor]{$1$ Is Not a Successor}
\end{itemize}
\end{dependencies}
